\section{Limitations and Failure Analysis}
\label{sec:limitations}

This section lists what is genuinely unresolved. Items that were found and fixed
are described in their own sections and are not repeated here as if they were
still open.

\paragraph{Free-stream preservation is not exact.} The uniform-flow check reaches
$O(10^{-11})$ rather than the requested $10^{-12}$, and the solver emits a warning
on every run. The cause is identified: the least-squares normal matrix
\cref{eq:normaleq} is poorly conditioned on the most stretched wall cells (aspect
ratios of order $10^3$), so the computed gradient of a constant field is not
exactly zero. The residual level involved is far below any converged residual, so
it does not affect the reported results, but it does place a floor on the
attainable residual and is the first thing that would need attention before
chasing deeper convergence. A weighted QR solve, or scaling the stencil offsets
before forming the normal equations, would improve it.

\paragraph{The CFL ramp can overshoot the stability limit and lose convergence.}
An earlier draft of this report attributed the transonic residual plateau to
limiter chatter in the smallest wall cells. \textbf{That explanation was wrong},
and the correction is instructive enough to record rather than quietly replace.

Examining the residual history of \texttt{naca0012\_m080\_inviscid} shows the
residual reaching $\num{4.6765e-4}$ at step 2546 with the ramped CFL at $49.7$ ---
that is $3.92$ orders, essentially the requested $4.00$ --- and then \emph{degrading}
to $\num{4.9518e-3}$ by step 3496 as the ramp pushed the CFL to its cap of $100$.
The run therefore discarded slightly more than one order of convergence and
reported the worse figure. The mechanism is not limiter chatter but CFL-driven
degradation past the stability limit of the scheme at this Mach number: the
geometric ramp of \cref{eq:ramp} continues to its cap regardless of whether the
iteration is still contracting.

The solver now retains the best state by residual norm and restores it if the march
ends more than a factor of two above that level, recording the fallback in the run
notes so that the reported residual and the written field always correspond. The
broader lesson is that a CFL cap taken from a case file is an input, not a
guarantee: the largest \emph{usable} CFL is a property of the scheme and the mesh,
and a ramp that ignores whether the residual is still falling can undo its own
progress.

With the fallback in place this case reports $3.92$ orders rather than $2.89$, so what
two earlier drafts described as a plateau an order short of target is in fact a near-miss
of the target caused by the ramp. \Cref{sec:naca} gives the full sequence. The limitation
that remains is real but narrower than first stated: the ramp is open-loop, and a CFL
schedule that responded to whether the residual was still falling would not need a
best-state fallback to protect it from itself.

\paragraph{A process failure in the run harness, recorded because it nearly corrupted
results.} During the re-runs required by the convergence fixes above, a termination
signal matched an outer wrapper script but orphaned the inner loop, so a subsequent
relaunch had two independent streams iterating the same seven cases into the same output
directories concurrently. The symptom was visible in the progress log --- cases completing
in an order the launching script does not use --- and the affected outputs were discarded
rather than harvested.

This is worth recording for two reasons. First, it is a reminder that results written
into a shared directory carry no record of which process wrote them, so a concurrent
writer produces files that are individually well-formed and collectively meaningless.
Second, it is the reason this report's harvesting script enforces a trusted-after
timestamp: every case whose \texttt{run\_status.json} predates the cutoff is reported as
still in progress rather than as a result, which is what turned a potential silent
corruption into a visible gap. The cutoff was raised five times during preparation as
successive fixes superseded earlier runs, and each raise was accompanied by re-running
every affected case with a single identical binary.

\paragraph{LU-SGS degrades at high CFL.} The spectral radius of the SGS iteration
approaches unity as $\mathrm{CFL}\to\infty$, so at the ramped targets of
\cref{tab:controls} the inner solve of \cref{eq:steadysystem} reaches its
iteration bound with the linear residual only partially reduced. The available
remedies are more inner sweeps (expensive) or a lower CFL cap (slower outer
convergence). What was deliberately \emph{not} done is inflating the diagonal to
make the reported inner ratio look better, which as noted in \cref{sec:lusgs}
produces a number that is meaningless. A matrix-free GMRES with the present
operator as preconditioner would be the natural improvement.

\paragraph{Boundary-layer resolution at $\Reinf=5000$.} The supplied NACA mesh was
not built for viscous computation at this Reynolds number. With
$\mu=\num{2e-4}$ the laminar boundary layer near the trailing edge has a thickness
of order $10^{-2}$ chords, spanned by only a few cells over much of the chord. The
skin-friction distributions of \cref{fig:naca-lam-cf} should therefore be read as
under-resolved, and the friction drag in \cref{tab:forces} carries a
correspondingly larger discretization error than the pressure drag. Mesh
refinement normal to the wall, not a solver change, is what this needs.

\paragraph{Shock capturing is monotone but diffuse.} The Venkatakrishnan limiter
smears a shock over two to three cells. At $\Minf=2.0$ the bow shock stands off
the leading edge in a region of relatively coarse cells, so the captured shock is
thicker there than on the body. This affects the sharpness of the contours in
\cref{fig:naca200inv-field} more than it affects the integrated forces, since the
pressure jump across the shock is captured conservatively even when it is spread.

\paragraph{The observed order of accuracy is not measured.} The evidence for
second order here is structural, diagnostic and comparative: exact recovery of
linear-field gradients ($\vLsqExactness$), the verified analytic Jacobian, the
positivity-fallback counts showing the second-order path is active in production,
and the first-versus-second-order comparison of \cref{sec:ordersens}, where first
order inflates cylinder drag by $60.5\,\%$ and out of the literature range. That
last result demonstrates the reconstruction does substantial work, but it is not
the same thing as measuring a convergence rate. A grid-convergence study on a
sequence of systematically refined meshes would establish the observed order
directly; only one mesh per geometry was supplied, so that measurement is absent.
This remains a gap in the verification and is better stated than glossed over.

\paragraph{Vorticity is a derived quantity.} The vorticity written to the field
file and plotted in \cref{fig:cyl200-wake} is computed from the same
least-squares velocity gradients used by the viscous flux, so it inherits their
accuracy. It is adequate for identifying and counting vortices, and hence for the
Strouhal estimate, but it is not a high-order vorticity reconstruction and its
peak values near the wall should not be read quantitatively.

\paragraph{Two-dimensionality of the $\Reinf=200$ wake.} The computation is
two-dimensional by construction. A real cylinder wake at $\Reinf=200$ is already
weakly three-dimensional, so the computed lift amplitude is expected to exceed the
experimental value even when the Strouhal number agrees well. This is a limitation
of the problem statement rather than of the solver, but it matters when comparing
the amplitude in \cref{sec:cyl200} against experiment.

\paragraph{The build locates more dependencies than it uses.} The CMake
configuration finds ParMETIS, Eigen and \texttt{fmt}, but no source file includes
any of them: the actual dependencies are CGNS, HDF5, METIS, zlib,
\texttt{nlohmann/json} and MPI. The macro \texttt{CNS2D\_HAVE\_EIGEN} is defined
but never used. This is stated because a build that advertises more than it needs
can mislead a reader about what the solver actually relies on, and because
partitioning is METIS serial $k$-way, not ParMETIS, despite ParMETIS being located.

\paragraph{The reproduce sequence points into \texttt{scratch/}.} The documented
reproduction path references \texttt{scratch/production.sh} and
\texttt{scratch/scaling\_study.sh}, which are the scripts that actually produced
the submitted runs. That directory is therefore not purely disposable. The
alternative would be retyping the commands into the README, where they could drift
from what was really executed; an honest pointer to the real script is preferred
over a tidier one that might not match.

\subsection{What would be done next}

In priority order: replace the least-squares normal-equation solve with a scaled
QR factorisation to recover exact free-stream preservation; add a matrix-free
GMRES outer Krylov solve preconditioned by the present LU-SGS operator, to
decouple outer convergence from the CFL cap; generate a refined mesh family to
measure the observed order of accuracy directly; and refine the airfoil meshes
normal to the wall so the $\Reinf=5000$ skin friction becomes mesh-independent.
One further experiment is specifically indicated by \cref{sec:floorsens}: repeat
the linear-wave floor sweep on the $\Minf=0.80$ and $\Minf=2.0$ cases, which is
the regime where the Roe carbuncle is actually notorious and the only place the
floor might prove necessary.

\section{Conclusions}
\label{sec:conclusions}

An original second-order finite-volume solver for the two-dimensional compressible
Navier--Stokes equations has been implemented, verified, and applied to all eight
required benchmark cases. The discretization is a cell-centred conservative finite
volume scheme with inverse-distance least-squares primitive reconstruction, a
Venkatakrishnan limiter, and a Roe flux whose Harten--Hyman entropy fix is extended
to the linear waves --- an extension that proved necessary, not decorative, because
the undamped entropy and shear waves at stagnation points otherwise stall the
steady residual and break symmetry. Steady cases use an implicit pseudo-time march
with local time stepping and a matrix-free LU-SGS inner solve; the vortex-shedding
case uses a true dual-time BDF2 scheme with frozen history levels and an outer
physical-time loop. Parallelism is METIS $k$-way partitioning of the cell dual
graph with neighbour-scoped halo exchange, with no full-mesh or full-state
replication during iterations.

The verification evidence is quantitative and is collected in
\cref{tab:verification}: geometric identities at machine precision, exact
least-squares gradients for linear fields, discrete conservation at
$\vConservation$, and an analytic Jacobian confirmed against finite differences.
Force coefficients agree to within $\num{5e-4}$ relative across $np=1,2,4,8$.

The most useful methodological result is negative and is documented in
\cref{sec:stationarity}: normalising the steady convergence test by the
impulsive-start residual can terminate a run while its forces are still moving,
producing results that look converged and are not. Requiring force stationarity as
well as residual reduction fixed this, and the affected cases were re-run. The
per-case status in \cref{tab:runstatus} is reported exactly as the solver recorded
it.
